% !TEX program = xelatex
\documentclass[UTF8,11pt,a4paper,fontset=windows]{ctexart}

\usepackage[a4paper,margin=1in]{geometry}
\usepackage{fontspec}
\setmainfont{Times New Roman}
\setsansfont{Arial}
\setmonofont{Consolas}
\usepackage{booktabs}
\usepackage{tabularx}
\usepackage{array}
\usepackage{longtable}
\usepackage{amsmath,amssymb,mathtools}
\usepackage{enumitem}
\usepackage{graphicx}
\usepackage{caption}
\usepackage[hidelinks]{hyperref}

\hypersetup{
  pdftitle={廉价磋商在狼人杀游戏中的应用},
  pdfauthor={Zhuo Chen; Jiaju Yang; Ruozhou Du}
}

\setlength{\parindent}{2em}
\setlength{\parskip}{0pt}
\setlength{\emergencystretch}{2em}
\raggedbottom
\setlist[itemize]{leftmargin=2.2em,itemsep=0.15em,topsep=0.25em}
\setlist[enumerate]{leftmargin=2.4em,itemsep=0.15em,topsep=0.25em}
\renewcommand{\arraystretch}{1.18}
\captionsetup{font=small,labelfont=bf}
\newcolumntype{Y}{>{\raggedright\arraybackslash}X}
\newcommand{\term}[1]{\textbf{#1}}

\title{\bfseries 廉价磋商在狼人杀游戏中的应用\\[0.4em]
\large 精确信念、结构化交流与模型条件终局收益}
\author{%
Zhuo Chen \quad Jiaju Yang \quad Ruozhou Du\\
\small 香港大学\\
\footnotesize\href{mailto:u3657972@connect.hku.hk}{u3657972@connect.hku.hk} \quad
\href{mailto:u3656084@connect.hku.hk}{u3656084@connect.hku.hk} \quad
\href{mailto:u3651313@connect.hku.hk}{u3651313@connect.hku.hk}\\[0.35em]
\small 指导教师：Jiayang Li 教授
(\href{mailto:jiayangl@hku.hk}{jiayangl@hku.hk})}
\date{2026年9月5日}

\begin{document}
\maketitle

\begin{abstract}
本中期研究考察狼人杀中的廉价磋商如何改变第一天公开票型，并据此构建动作级收益矩阵。理论基线是 Shitong Wang 的《Optimal Strategy in the Werewolf Game: A Theoretical Study》，该研究分析较简单的角色组合与诚信约束下的预言家信息公开。团队将其扩展为包含两名狼人、两名平民、一名女巫、一名预言家和一名守卫的固定七人场景，用于研究顺序发言和角色能力交互条件下策略交流的演化。报告汇报两项工作：其一是重设计游戏引擎，使目标规则、智能体顺序发言、角色视角动态上下文和实验处理提示词保持一致；其二是构建精确信念 Cheap-talk 决策矩阵，对合法隐藏站位进行精确条件化，并以固定功利等级策略下的蒙特卡洛模拟推演终局。一个生产请求在三档证据强度下评价22个预言家动作，每格使用100条轨迹；另有16局语言模型对局用于提供首日票型机制的探索性证据。当前结果支持研究工具具有可复现性和信息隔离，但不能据此断言扩展博弈均衡、普适最优发言策略或已经校准的行为模型。
\end{abstract}

\section{概述}

\subsection{研究目的与范围}

廉价磋商（cheap talk）是指不直接改变游戏正式规则、没有强制约束力且本身不产生成本，但被接收和解释后可能改变信念与行动的交流。狼人杀中的隐藏身份和非对称信息使公开发言紧接集体投票发生，适合对该机制进行受控研究。本中期阶段的主要目标是识别结构化发言如何改变第一天公开票型，并通过收益矩阵量化这种变化最终对应的阵营终局效用。

报告包含两项相互衔接的研究产出。第一项是实验游戏引擎的重设计：规则、提示词、发言顺序和动态上下文经过统一，使观测到的第一天票型只能由每次行动发生前可用的信息产生。第二项是\term{精确信念 Cheap-talk 决策矩阵}：信念层依据当前行动者的合法信息精确条件化有限隐藏站位，决策层从完整后验分布抽样，并在显式功利等级策略下把每个候选声明推演至终局效用。自然语言模型用于生成探索性完整对局记录，但不生成矩阵内的后续行动，也不决定矩阵排名。

固定板子包括两名狼人、两名平民、一名女巫、一名预言家和一名守卫。相较 Wang 的简化理论基线，该组合有意同时引入能够协作的知情少数、缺少身份能力的发言者、可验证的私有信息以及两种额外夜间能力。其规模仍允许对当前隐藏角色进行精确枚举，但发言后果不再由单一的信息公开决策决定。因此，该板子构成从可处理理论模型通向复杂社会推理场景的实验桥梁，用于研究顺序上下文、欺骗、保护与有限资源共同作用时廉价磋商的演化。

第一天的主要干预包括平民虚假声明预言家身份以分担风险、真实预言家不公开信息，以及狼人虚假声明预言家身份。普通指控、支持、投票意向、中性对照动作和主动沉默也进入候选集合。矩阵保留每个具体声明的目标及声称的查验结果。首日票数集中程度、被放逐席位的身份类别、身份声明竞争和预言家存活构成连接发言与终局收益的可观测机制，但不能替代收益本身。

两类证据保持分离。一个生产预算的预言家请求为22个动作、3档证据强度提供模型条件收益，每格使用100条轨迹。较早语料包含16局分析用语言模型对局和1局冒烟运行，可说明战术是否显现及首日票型如何变化，但不能校准矩阵或支持总体结论。

\subsection{研究目标}

本研究设置四项目标：

\begin{enumerate}
  \item 将 Wang 受诚信约束的信息公开模型扩展至顺序化且可能包含欺骗的交流，同时不把其均衡结论外推至不同博弈。
  \item 对当前隐藏角色保持精确推理，并以可控成本评价组合规模迅速增长的未来。
  \item 在匹配随机条件下，估计候选发言相对中性参照对首日票型机制与终局效用的影响。
  \item 区分可复现的矩阵证据与语言模型探索性对局所提供的有限机制证据。
\end{enumerate}

\subsection{核心术语}

\term{行动者视角}指公开状态与规则依法向当前行动者揭示的私有事实之和。每名玩家知道自己的角色；狼人知道狼队友；预言家知道已完成查验的阵营结果；普通死亡不会自动公开角色。\term{隐藏世界}是与固定角色多重集合相容的一种完整站位。\term{后验分布}是在行动者合法信息与结构化公开证据条件下，对这些隐藏世界形成的概率分布。\term{终局轨迹}是从当前发言决策开始延伸至好人或狼人阵营胜利的一次模拟续局。

\term{矩阵对照动作}表示沟通强度为零、且不包含指控、支持、投票意向、身份声明、查验结果或战术标签的中性结构化发言轮次事件。它不同于\term{显式沉默}；后者是作为独立候选进入排序的主动策略动作。二者也不同于\term{整局实验对照组}；该对照组不启用平民或狼人虚假声明预言家身份的处理，同时要求真实预言家提供有效发言。只有保持这三个概念的区分，比较结果才具有清晰语义。

\section{发现问题：对 Wang 理论基线的演化拓展}

\subsection{理论基线研究的既有结论}

Shitong Wang 的《Optimal Strategy in the Werewolf Game: A Theoretical Study》构成本项目的理论基线\cite{wang2025}。对于没有预言家的对局，该研究质疑传统随机投票假设，并提出“random strategy+”：当知情少数能够制造平票时，加入协调一致的全押响应。在其规则下，该策略对狼人阵营弱支配传统随机策略，夜间自我淘汰也被证明为严格劣势。对于含预言家且采用诚信规则的对局，Wang 把隐藏或公开信息建模为预言家信息集上的递归贝叶斯决策，并推导出能够诱导完美贝叶斯均衡的策略组合。

Wang 还把取消诚信规则、引入顺序发言以及加入守卫等额外角色列为开放方向。这些扩展会实质改变博弈：虚假公开声明取消了预言家发言的规则性可信度；顺序发言使早期声明进入后续玩家的信息集；女巫与守卫使夜间存活受到私有能力和有限资源约束。Wang 的递归公开结论为本研究提供动机，但不能原样移植到该复杂设定。

两项研究之间的关系如下表所示。

\begin{table}[htbp]
\centering
\caption{继承理论与团队扩展的边界}
\begin{tabularx}{\textwidth}{>{\bfseries}p{0.19\textwidth}YYY}
\toprule
维度 & Wang 的理论基线研究 & 继承的假设 & 团队在中期阶段的扩展 \\
\midrule
信息 & 隐藏阵营与预言家查验 & 决策必须遵守非对称信息边界 & 对完整合法站位执行逐行动者精确条件化 \\
交流 & 同时公告；预言家分析采用诚信规则 & 公开信息可能影响后续行动 & 顺序化结构声明，包括真实与虚假信息 \\
策略对象 & 预言家隐藏或公开的选择 & 公开时机具有终局后果 & 面向每名发言者的“具体动作×目标”收益矩阵 \\
角色 & 平民、狼人、预言家 & 两个对立阵营及阵营一致的终局效用 & 在固定七人规则中加入女巫与守卫 \\
求解概念 & 在特定约束下递归求最优动作与 PBE & 信念与续局价值必须显式建模 & 模型条件蒙特卡洛价值；不作均衡声明 \\
开放方向 & 欺骗、顺序与额外角色 & 上述缺口具有理论意义 & 同时把三项因素显式化并使其可观察 \\
\bottomrule
\end{tabularx}
\end{table}

\subsection{为何不受诚信约束的廉价磋商改变了问题}

在诚信规则下，公开的预言家身份与查验可以按规则被视为可信。取消该规则后，相同语义形式可能来自三个具有不同策略动机的来源：真实预言家公开合法私有信息、平民试图转移夜间攻击，或狼人制造身份声明。接收者必须同时更新对目标与发言者的信念。虚假声明可能通过改变投票或攻击目标帮助发言者阵营，但也可能引发反噬、排挤真实信息，或造成后续声明之间的不一致。因此，发言价值取决于谁发言、行动者知道什么、处于何种发言位置、指向哪个目标，以及后续玩家把该声明当作多强的证据。

第一天研究范围提供了清晰干预点。每名存活玩家按固定顺序获得一次公开发言机会；后续发言者只看到此前已发生的结构化事件；随后所有具备投票权的玩家公开投票。在目标规则集中，最高票平票通过在并列玩家中均匀随机抽取一人解决；死亡仅公开席位而不公开角色。这些规则保留了社会推理问题的核心：发言改变公开证据，却不会机械地改变发言者真实身份。

隐藏世界有限但规模不可忽略。七个有标号席位分配两名狼人、两名平民、一名女巫、一名预言家和一名守卫时，完整角色站位数为

\[
|\Omega|=\frac{7!}{2!\,2!}=1260.
\]

\begin{description}[style=nextline,leftmargin=2.9cm,labelwidth=2.5cm]
  \item[$\Omega$] 七个有标号席位全部完整合法角色站位组成的有限集合。
  \item[$|\Omega|$] 该集合的元素数量；同类狼人和平民分别可交换时共有1260种。
  \item[$7!$] 暂时把七个角色标记视为彼此不同所得到的排列数。
  \item[$2!\,2!$] 分别消除两名同类狼人和两名同类平民重复计数的修正项。
\end{description}

在当前决策点执行精确枚举是可行的，但穷举全部未来组合并不可行。顺序声明、逐人投票、平票、狼刀、查验、药物、守护、死亡和后续轮次会使续局数量相乘。此外，语法上合法的分支数量不是行为概率；对所有分支等权会把组合结构误当作行动选择模型。

\subsection{可证伪的研究缺口}

研究缺口是缺少一个能够在角色视角约束下评价欺骗或不公开信息的稳定估计对象。所需反事实对照保持行动者状态与随机条件不变，用一个具体结构化声明替换当前中性事件，并测量终局阵营效用的变化。估计过程还必须区分有限抽样误差与模型不确定性。

本项目排除四类捷径：观察者独有的完整站位不得进入行动者请求；语言模型不提供未来行动概率；完整树路径数不代表行为概率；面向不同目标的虚假预言家声明不得合并为单一分数。这些约束用于保持可复现性与可解释性。

由此提出的假说都具有条件性：在固定模型下，非零发言证据应使部分动作相对矩阵对照动作发生价值变化；具有不同目标的声明应出现收益异质性；行动者无法区分的角色视角应生成相同请求；配对随机流应使候选动作减去矩阵对照动作的估计不依赖执行顺序。语言模型对局提供另一类机制检查：启用的战术应在文本中显现并可能改变首日票型，但小样本分组不具备推断一般胜率的统计效力。

\section{提出解决方案}

\subsection{两层决策模型}

所提出的方法把当前信息不确定性与未来状态空间的组合增长分开处理。在决策时刻，信念层枚举与固定板子和行动者硬知识相容的全部隐藏世界，再纳入结构化公开证据并归一化剩余权重。行动者后验分布定义为

\[
B_t^{(i)}(\theta)=
\frac{B_{t-1}^{(i)}(\theta)L_t^{(i)}(\theta)}
{\sum_{\theta'}B_{t-1}^{(i)}(\theta')L_t^{(i)}(\theta')}.
\]

\begin{description}[style=nextline,leftmargin=3.4cm,labelwidth=3.0cm]
  \item[$B_t^{(i)}(\theta)$] 行动者 $i$ 在决策时刻 $t$ 对隐藏站位 $\theta$ 赋予的后验概率。
  \item[$B_{t-1}^{(i)}(\theta)$] 纳入当前证据之前，行动者对同一站位的概率判断。
  \item[$L_t^{(i)}(\theta)$] 在站位 $\theta$ 下出现新增结构化证据的似然，并从行动者 $i$ 的合法信息视角计算。
  \item[$\theta$] 一个完整合法隐藏角色站位；$\theta'$ 是对全部相容站位求和时使用的索引。
  \item[$i,t$] 分别表示当前行动者与顺序证据更新索引。
\end{description}

该层不使用蒙特卡洛近似，也不使用 top-$k$ 截断。如果证据使全部相容世界权重为零，系统报告模型不一致，而不是静默恢复均匀分布。硬知识也采用严格定义：自身角色、狼人依法知道的队友信息、预言家已经完成的阵营查验，以及规则明确公开的真实身份事件可以进入；公开声明、未公开身份的死亡、怀疑或高后验概率不能升级为确定身份知识。

决策层对当前发言事件实施干预，并把对局继续到终局。具体动作 $a$ 在证据强度 $c$ 下的理论价值为

\[
V^{(i)}(a,c)=
\mathbb{E}_{\Theta\sim B_t^{(i)},\;
\tau\sim P_{\rho,c}(\cdot\mid \operatorname{do}(a),\Theta,I_t^{(i)})}
\left[u_i(\tau)\right].
\]

\begin{description}[style=nextline,leftmargin=3.4cm,labelwidth=3.0cm]
  \item[$V^{(i)}(a,c)$] 行动者 $i$ 在证据强度 $c$ 下执行动作 $a$ 的理论模型条件期望终局效用，即目标估计量。
  \item[$\mathbb{E}$] 同时对隐藏世界不确定性与随机续局轨迹求期望的算子。
  \item[$\Theta$] 直接从行动者完整后验分布抽取的随机隐藏站位。
  \item[$\tau$] 从当前干预延伸至终局的一条随机续局轨迹。
  \item[$P_{\rho,c}$] 由后续策略 $\rho$、证据强度 $c$ 和显式规则随机性诱导的轨迹分布。
  \item[$\operatorname{do}(a)$] 在保持相同初始状态时，以动作 $a$ 替换当前中性发言事件的干预。
  \item[$I_t^{(i)}$] 时刻 $t$ 行动者 $i$ 依法可获得的公开与私有信息。
  \item[$u_i(\tau)$] 阵营一致的终局效用：行动者阵营胜利取 $+1$，否则取 $-1$。
\end{description}

该定义明确了证据边界。研究观察者可以保留完整角色站位用于事后审计，但该站位不进入行动者收益查询或目标排序。因此，矩阵是事前决策工具，而不是事后上帝视角评分。

\subsection{结构化动作与三种不同对照}

自然语言具体措辞被有意排除在动作键之外。一级动作族包括矩阵对照动作、指控、支持、投票意向、预言家身份声明和显式沉默。实际评价发生在二级动作层；每一行固定相关目标、声称角色、声称的查验目标与阵营结果、沟通强度和实验战术标签。完整虚假预言家声明必须指向另一名存活玩家，并声称该目标为好人或狼人；不给查验的弱身份声明属于独立行；真实预言家只能公开实际私有查验；主动沉默以空内容推进发言位置。

\begin{table}[htbp]
\centering
\caption{不得混淆的三类对照概念}
\begin{tabularx}{\textwidth}{>{\bfseries}p{0.22\textwidth}YY}
\toprule
概念 & 操作定义 & 统计或实验用途 \\
\midrule
矩阵对照动作 & 经过预定发言轮次，不加入战术内容，沟通强度为零 & 用于计算单一状态内的配对动作差 \\
显式沉默 & 主动不提供公开内容的独立候选动作 & 把不公开信息作为策略排序；当前可能与矩阵对照动作数值相同 \\
整局实验对照组 & 重复对局中不启用平民或狼人虚假预言家声明，真实预言家必须有效发言 & 在完整对局之间比较处理组与对照组 \\
\bottomrule
\end{tabularx}
\end{table}

即使两个动作当前得到相同数值，也必须保留概念差异。矩阵对照动作表示缺少待评价干预；显式沉默表示经选择的省略行为；整局实验对照改变整场对局允许的战术集合，不能替代前两种动作级对象。

\subsection{证据强度与后续行为}

接收者对结构化声明的敏感程度由证据强度参数表示。本研究固定评价0、0.5和0.8三档。这些值不是“声明为真”的概率，而是在中性单位似然与指定发言似然之间进行插值：

\[
L_c(e\mid\theta)=(1-c)+cL_{\mathrm{speech}}(e\mid\theta),
\qquad c\in\{0,0.5,0.8\}.
\]

\begin{description}[style=nextline,leftmargin=3.4cm,labelwidth=3.0cm]
  \item[$e$] 作为公开证据纳入更新的结构化发言事件。
  \item[$L_{\mathrm{speech}}(e\mid\theta)$] 隐藏站位 $\theta$ 下的基础发言似然。
  \item[$L_c(e\mid\theta)$] 证据强度为 $c$ 时的插值似然。
  \item[$c$] 取值位于 $[0,1]$ 的敏感性分析参数；本报告分别保留0、0.5和0.8三档。
  \item[$1-c$] 中性似然的权重；当 $c=0$ 时，声明不提供身份识别证据。
\end{description}

后续模拟玩家采用固定功利等级策略。游戏硬规则优先于实验战术约束，后者再优先于阵营目标、身份能力优先级和同级动作等权。五档规范分数1、0.75、0.5、0.25和0通过 softmax 函数转换为概率分布：

\[
\pi_{\rho}(a\mid s,i)=
\frac{\exp(Q_i(s,a)/\tau_{\mathrm{policy}})}
{\sum_{a'\in\mathcal A_i(s)}\exp(Q_i(s,a')/\tau_{\mathrm{policy}})},
\qquad \tau_{\mathrm{policy}}=0.25.
\]

\begin{description}[style=nextline,leftmargin=3.4cm,labelwidth=3.0cm]
  \item[$\pi_{\rho}(a\mid s,i)$] 策略 $\rho$ 在模拟状态 $s$ 中为行动者 $i$ 选择合法动作 $a$ 的概率。
  \item[$Q_i(s,a)$] 动作的规范化功利等级分数，只使用上述五档。
  \item[$\mathcal A_i(s)$] 状态 $s$ 中行动者 $i$ 的完整合法动作集合；$a'$ 是该集合上的求和索引。
  \item[$\tau_{\mathrm{policy}}$] 固定为0.25的 softmax 温度，控制不同动作等级之间的概率离散程度，并与证据强度 $c$ 相互独立。
  \item[$\exp$] 把等级分数转换为正的归一化权重所使用的指数函数。
\end{description}

该策略透明且可复现，但属于建模假设。它不从历史对局估计，不由语言模型提供，也不被宣称为最优或均衡行为。未来模拟行动者不会递归构建自己的矩阵。该非递归约束保持目标分布固定，并避免决策问题嵌套造成不可控增长。

\subsection{终局收益估计与配对比较}

生产规范为每个具体动作、每个证据档位执行 $N=100$ 条终局轨迹。普通蒙特卡洛估计及其标准误为

\[
\widehat V_N^{(i)}(a,c)=\frac{1}{N}\sum_{n=1}^{N}u_i(\tau_n^{a,c}),
\qquad
\operatorname{SE}_N(a,c)=\frac{s_N(a,c)}{\sqrt N}.
\]

\begin{description}[style=nextline,leftmargin=3.4cm,labelwidth=3.0cm]
  \item[$N$] 单个“动作×证据”单元格的终局轨迹数；生产规范使用100，探索性矩阵案例使用10。
  \item[$n$] 从1至 $N$ 的样本索引。
  \item[$\tau_n^{a,c}$] 动作 $a$、证据强度 $c$ 和样本索引 $n$ 对应的终局轨迹。
  \item[$\widehat V_N^{(i)}(a,c)$] 估计理论值 $V^{(i)}(a,c)$ 的样本均值。
  \item[$s_N(a,c)$] 已实现终局效用的样本标准差。
  \item[$\operatorname{SE}_N(a,c)$] 由有限模拟轨迹数量造成的蒙特卡洛标准误。
\end{description}

在固定证据档位和样本索引下，全部候选复用隐藏世界抽样、后续行动随机流、平票随机数和其他规则随机性。令 $a_0$ 表示矩阵对照动作，则配对差及其标准误为

\[
\widehat\Delta_N^{(i)}(a,c)
=\frac{1}{N}\sum_{n=1}^{N}
\left[u_i(\tau_n^{a,c})-u_i(\tau_n^{a_0,c})\right],
\]

\[
d_n(a,c)=u_i(\tau_n^{a,c})-u_i(\tau_n^{a_0,c}),
\qquad
\operatorname{SE}_{\Delta,N}(a,c)
=\frac{1}{\sqrt N}
\sqrt{\frac{1}{N-1}\sum_{n=1}^{N}
\left[d_n(a,c)-\widehat\Delta_N^{(i)}(a,c)\right]^2}.
\]

\begin{description}[style=nextline,leftmargin=3.4cm,labelwidth=3.0cm]
  \item[$a_0$] 在相同行动者、状态、证据强度和策略条件下的中性矩阵对照动作。
  \item[$\widehat\Delta_N^{(i)}(a,c)$] 样本内“候选减对照”的终局效用差均值。
  \item[$d_n(a,c)$] 样本索引 $n$ 对应的配对差。
  \item[$\tau_n^{a_0,c}$] 与候选轨迹共享隐藏世界和具名随机流的对照轨迹。
  \item[$\operatorname{SE}_{\Delta,N}(a,c)$] 配对差均值的标准误，样本方差使用 $N-1$ 个自由度计算。
\end{description}

共同随机数能够降低与动作无关的随机差异，使干预比较更容易解释；它不能消除模型设定错误，也不能把条件模拟差转换为现实因果效应。

每个单元格以向量呈现，而不是聚合为单一策略分数：

\[
\mathcal M^{(i)}_{a,c}=
\left(\widehat V_N^{(i)},\operatorname{SE}_N,
\widehat\Delta_N^{(i)},\operatorname{SE}_{\Delta,N},
N,\mathbf n_{\mathrm{scenario}}\right)_{a,c}.
\]

\begin{description}[style=nextline,leftmargin=3.4cm,labelwidth=3.0cm]
  \item[$\mathcal M^{(i)}_{a,c}$] 行动者 $i$、动作 $a$ 与证据强度 $c$ 所共同确定的完整报告单元格。
  \item[$\mathbf n_{\mathrm{scenario}}$] 六个互斥反应情景的计数向量：行动者反噬、目标成功转移、声明被接受、声明受到争议、发言被忽略和其他。
  \item[其他分量] 前文刚刚定义的估计量、标准误与样本数；下标 $a,c$ 表示所有分量均对应同一单元格。
\end{description}

轨迹生成后按照固定优先级分配情景标签，默认信念变化阈值为0.05。情景只用于解释模拟机制，不反向影响收益或策略。

\section{行动以及验证猜想}

团队推进了两条相互衔接的工作线：重设计实验游戏引擎，并构建收益矩阵估计器。引擎工作使受控的第一天票型实验成为可能，矩阵工作则把该研究问题转换为可复现的动作级估计量。

\subsection{游戏引擎重设计}

引擎围绕固定七人板子和显式时间顺序重新配置。夜间先由狼人私下协商并投票确定攻击目标，再由预言家、女巫和守卫依次行动，最后统一结算。白天由存活智能体按固定顺序各发言一次，再提交公开放逐票；平票均匀随机解决，同时排除警徽、遗言和死亡公开身份。女巫资源、守卫限制、狼人私下协作以及可选特殊攻击均由明确规则表达，而不只依赖提示词约定。该工作避免未记录房规改变首日票型机制。

提示词同时在稳定规则层和实验处理层得到改进。每名智能体接收其合法角色、能力、剩余资源、当前可执行动作以及所属实验条件启用的战术。处理指令置于 system 层，并一致作用于发言、投票和夜间行动。整局对照条件移除虚假预言家处理，并要求真实预言家交流合法查验，使较低优先级的对话指令不能静默破坏处理隔离。

动态上下文按照事件时间和信息可见性修正。发言者只接收此前的公开声明、公开票型、死亡信息及合法私有事实；后置位声明不能进入前置位提示，狼人协商仅对本阵营可见，预言家查验在公开前保持私有。完整发言、票型、私有理由和终局结果被保留用于审计。重设计协同修改规则、阶段顺序、动作验证、可见性、提示词组装与文本记录。早期引擎测试集共有133项通过，定向冒烟检查覆盖夜序、狼人私聊、动作合法性、文本渲染和完整对局运行。这些检查体现工程投入，而非本报告的主要科学结果。

\subsection{收益矩阵构建}

行动者视角边界被转化为可检验条件。矩阵请求标识行动者、行动者自知角色、依法知道的队友或查验事实、公开状态、精确后验摘要、完整候选集合、证据档位、样本目标和种子方案。属于研究观察者的完整站位会在决策通道被拒绝。因此，只要行动者观察相同，即使观察者知道不可见角色标签发生变化，请求仍应保持一致。

发言被转换为可审计干预。每个候选在模拟前生成，并具有稳定语义身份；全部同级目标显式保留；真实预言家声明受到实际私有查验约束；完整虚假声明必须包含目标和阵营结果；主动沉默与中性参照保持分离。该表示使同一动作词汇能够描述模拟矩阵与经过编码的自然语言对局记录。

前向估计采用匹配抽样。一个批次对相同证据档位和样本索引区间内的全部候选同时评价，只保留聚合矩、配对差及情景计数。完成批次采用幂等提交，中断运行以相同种子恢复缺失索引，失败运行不作为完整排名呈现。在每条轨迹中，第一天票型是直接观测的机制，阵营胜负则提供终局效用。因此，矩阵把票型变化继续传播至剩余对局，而不把首日放逐结果本身等同于收益。

\subsection{假说与验证标准}

当前验证围绕五项假说进行。这些是假设条件下的实现与模型行为假说，不是关于未知人类总体的结论。

\begin{table}[htbp]
\centering
\caption{中期假说与证伪标准}
\begin{tabularx}{\textwidth}{>{\bfseries}p{0.08\textwidth}YY}
\toprule
编号 & 条件性假说 & 能够证伪或削弱假说的证据 \\
\midrule
H1 & 角色视角条件化能够阻止观察者信息影响行动者矩阵。 & 在保持行动者视角不变时，改变行动者不可见角色会改变请求或收益计算。 \\
H2 & 共同随机数能够生成不依赖执行顺序的配对估计。 & 相同请求在1、2、4个工作进程、不同候选顺序或完成顺序下得到不同数值。 \\
H3 & 在指定模型下，结构化廉价磋商具有随目标与证据变化的价值。 & 在发言能够改变信念的状态中，全部非对照动作仍在所有目标和三档证据下完全相同。 \\
H4 & 矩阵完整性能够被机械审计。 & 缺少动作、证据列、样本、配对对照或情景计数的结果仍被标记为完整。 \\
H5 & 被启用的战术能够在完整语言模型对局中显现，并可能改变首日投票机制。 & 处理组文本没有出现预定战术，或公开动作违背处理配置却未作为偏离记录。 \\
\bottomrule
\end{tabularx}
\end{table}

H1 通过角色视角隔离检查和输入拒绝验证；H2 通过分别使用1、2、4个工作进程重复相同小型请求并逐值比较验证；H3 通过一个生产预算矩阵和十个探索性行动者视角案例观察；H4 通过样本计数闭合、反应计数闭合、确定性终局案例、持久化、查询、中断与恢复检查验证；H5 通过完整文本与首日投票结果在较早的语言模型对局语料中检查。

\subsection{矩阵评价方案}

生产参照采用3档证据强度，每个具体“动作×证据”单元格使用100条轨迹。它不同于用于低成本检查不同角色和目标的十视角探索性研究，后者每个单元格仅使用10条轨迹，共包含308个动作行和924个单元格。后者报告的标准误最高约为0.33，因此其数值案例只用于说明矩阵读取和敏感性，不用于建立稳定动作排名。

可靠性检查以完整单元格为粒度比较。确定性终局校验要求：当行动者阵营在任意续局中均已获胜时，效用必须为 $+1$，标准误必须为0。可复现性要求同一请求在不同工作进程数量下结果完全一致。完整性要求每个动作具备全部3档证据列、每格达到指定样本数，且六类情景计数之和严格等于该样本数。输入隔离要求拒绝上帝视角完整站位、矛盾私有事实和非法动作。在中期检查点，共有65项覆盖这些性质的模块级测试通过。

测试数量只说明研究工具的内部实现，不说明行为策略是否充分。一个完全可复现但设定错误的模型仍然设定错误。因此，数值正确性、规则正确性、信息隔离和外部行为验证被分别报告。

\subsection{语言模型探索性实验方案}

较早的廉价磋商实验在同质、理性、功利的人格配置下使用语言模型作为对局行动机制，共形成17次完整运行：10局分析用整局对照、1局额外对照冒烟运行，以及三个处理组各2局。主要分析排除冒烟运行，因此包含16局。对照组禁用平民和狼人的虚假预言家战术，并要求真实预言家如实交流；三个处理组分别启用狼人虚假声明预言家、平民虚假声明预言家或预言家不公开信息。狼人虚假声明组还开放两种夜间战术，但模型在观测运行中没有选择它们。

该语料的价值在于机制追踪：公开发言、投票、私有理由和终局均被完整记录。其统计局限同样明确：每个主要处理组只有2局；角色站位、发言顺序和语言模型抽样可能主导组均值；该语料还遵循采集时记录的实验引擎配置。它不与当前精确信念矩阵进行数值合并，因为当前矩阵冻结的规则集和行为策略定义了不同估计对象。因此，两类证据只进行定性机制对照：矩阵中的信念扭曲、票型转移、反噬和真实信息损失是否也会出现在机器玩家完整对局叙事中。

\section{观测结果}

\subsection{第一天票型观测}

主要的16局分析包括10局对照和6局处理对局。下表复现观测结果，用于机制讨论而非统计推断。

\begin{table}[htbp]
\centering
\caption{语言模型探索性对局结果}
{\small\setlength{\tabcolsep}{3pt}
\begin{tabularx}{\textwidth}{>{\bfseries\raggedright\arraybackslash}p{0.15\textwidth}*{5}{>{\centering\arraybackslash}p{0.075\textwidth}}Y}
\toprule
组别 & 局数 & 好人胜 & 好人胜率 & 平均天数 & 首日放逐狼人 & 主要观测 \\
\midrule
对照组 & 10 & 4 & 40\% & 2.3 & 50\% & 无虚假预言家声明；真实预言家报告查验。 \\
狼人虚假声明 & 2 & 1 & 50\% & 2.5 & 50\% & 两局均出现虚假声明；一局形成三方预言家声明竞争。 \\
平民虚假声明 & 2 & 0 & 0\% & 2.0 & 0\% & 两局首日均放逐真实预言家。 \\
预言家不公开 & 2 & 0 & 0\% & 2.0 & 0\% & 全局没有预言家身份声明；两局首日均放逐平民。 \\
\bottomrule
\end{tabularx}
}
\end{table}

在10局分析用对照中，没有非预言家玩家声称该角色，说明处理隔离在文本层面可见；真实预言家也按要求在首日交流。好人赢得4局，5局首日放逐对象为狼人。这些比例只描述当前观测样本。

两局狼人处理对局都出现虚假预言家声明。其中一局里，真实预言家报告好人查验，一名处于前置位的狼人向狼队友发出好人查验，狼队友随后以虚假狼人查验指控真实预言家。由此形成的三方身份声明竞争显示了协调保护与叙事竞争。好人在该局失败，但在第二局通过放逐两名狼人逆转。该对照说明战术确实显现且目标会影响互动，不能据此建立稳定的50\%处理胜率。

两局平民虚假声明对局呈现出最清晰的不利机制。每局都有平民声称预言家，同时真实预言家也发言；两局首日投票都放逐了真实预言家。虚假声明并未被观测为成功吸引夜间攻击，反而让多个身份声明污染公共信号，并把怀疑转向真实信息源。这一模式与矩阵中的“行动者反噬”和“声明受到争议”情景相容，但两例无法建立该机制的发生频率或期望收益。

两局预言家不公开对局中，预言家成功避免身份声明，但好人阵营失去主要公开信息来源；两局首日都放逐平民，狼人阵营很快达到人数优势。该观测说明生存与信息之间存在权衡：隐藏身份可能在短期保全角色，却降低阵营协调能力。其中一名预言家存活至终局，因此个体存活本身并不是充分的效用指标。

在两局狼人虚假声明对局的四个可观测夜晚中，空刀和自刀虽然可选，却从未被选择。该事实只表示这些动作在特定语言模型提示和局势下未被采用，并不能说明它们在当前扩展规则下为严格劣势；Wang 关于自我淘汰的结论适用于理论基线所采用的更简化假设。

\subsection{收益矩阵结果}

一个采用生产预算的七人预言家请求评价了22个具体动作和3档证据强度，共得到66个单元格，每格包含100条终局轨迹。终局效用估计范围为 $-0.58$ 至 $-0.34$。该范围只描述一个行动者可见状态和一个策略规范，不是跨场景的预言家通用胜率区间。负值表示在该特定请求所定义的模拟分布中，行动者阵营失败次数多于胜利次数。

十视角探索性报告提供了下表所示的多样动作案例。表内所有案例每格仅使用10条轨迹，而不是生产预算；标准误最高约0.33，且终局效用只取正负1，因此在该样本规模下估计较为粗糙。

\begin{table}[htbp]
\centering
\caption{探索预算下的说明性矩阵观测}
\begin{tabularx}{\textwidth}{>{\bfseries}p{0.24\textwidth}p{0.23\textwidth}Y}
\toprule
条件 & 记录数值 & 允许作出的解释 \\
\midrule
零证据档 & 大多数动作与各自中性参照相同 & 当发言作为身份信息被忽略时，这些案例中的信念通道贡献较小。 \\
平民虚假预言家声明，高证据档 & 对照 $-0.4$；特定声明 $+0.4$；配对差 $+0.8$ & 在该状态中，以该目标和声明结果替换中性事件后，模型条件估计由负变正。 \\
狼人虚假预言家声明，高证据档 & 对照 $+0.6$；特定声明 $+1.0$；配对差 $+0.4$ & 目标身份具有影响；包括部分涉及已知狼队友的其他声明可能具有较低收益。 \\
预言家持有狼人查验 & 中、高证据档的配对差均为 $+0.4$ & 在该信息状态中，如实公开得到正向配对估计；对应的好人查验案例没有相同差异。 \\
\bottomrule
\end{tabularx}
\end{table}

输出直接支持三项事实。第一，在没有额外且得到论证的加权分布时，3档证据列不能合并。第二，战术族过于粗糙，不能被赋予单一收益，因为目标和声称阵营会改变估计。第三，提高证据强度并不会单调提高发言者效用。更强的可信反应可能加强有益转移，也可能放大反噬或使有害欺骗产生更大后果。

模型允许一种条件性解释：在指定规则和策略下，结构化发言能够通过信念与后续行动通道改变期望终局效用。模型不允许把最高收益行解释为通用策略、支配动作或均衡收益；使用10条轨迹的探索性数值尤其容易受抽样波动影响。

\subsection{模型结果与观测机制的综合}

矩阵与语言模型语料在“可能机制”层面相互呼应。两者都显示声明的语义目标具有重要性，证据被更强地接收并不总是对发言者有利，隐藏或伪造信息也可能对本阵营形成负外部性。矩阵增加了反事实控制：多个动作在匹配隐藏世界和随机流下比较。语言模型语料增加了叙事真实性：当自然语言行动者进行完整对局时，竞争声明、怀疑和票型转移确实可能出现。

两类证据尚未在“校准后的效应大小”层面一致。矩阵使用人工指定的后续策略和3个敏感性锚点；语料只使用一种模型配置且处理组极小。矩阵输出的是某个决策状态条件下的动作级期望效用，而对局表报告的是整局处理分配下实现的终局。矩阵配对差 $+0.8$ 与某处理组观测胜率0\%并不矛盾，因为二者的状态、单位、抽样方案和估计对象均不同。

因此，当前结果对研究工具的验证强于对任何具体战术的验证。角色视角精确条件化、跨进程数量的确定性复现、终局解析校验、完整单元格审计和机制可见的对局文本支持研究继续推进；这些数据尚不能验证总体行为模型，也不能形成公开、欺骗或沉默的通用建议。

\section{结论}

Wang 的《Optimal Strategy in the Werewolf Game: A Theoretical Study》说明，在显式假设下，狼人杀中的信息结构与公开时机可以接受严谨建模。本研究沿其开放方向推进欺骗性交流、顺序发言和额外角色。团队首先重设计游戏引擎，使提示词、上下文可见性、行动顺序和规则能够支持受控的首日票型观测；随后对当前隐藏世界保持精确推理，并用受控、非递归的终局模拟替代未来完整展开。由此形成的矩阵从行动者合法视角评价具体发言动作，并报告终局效用、蒙特卡洛不确定性、相对中性参照的配对差和解释性反应计数。

本阶段的主要成果属于方法和有边界的证据。一个生产请求已经展示每格100条轨迹的规范流程；探索性矩阵显示收益随目标和证据强度变化；少量语言模型对局显示信息污染、身份声明竞争、票型转移与隐藏信息成本能够在完整对局中出现。可靠性检查支持研究工具内部的可复现性与信息隔离。

本研究没有求得扩展博弈的均衡。功利等级策略是显式续局投影模型，而不是每名玩家的已求解最优反应；证据强度是敏感性参数，不是经过校准的真实概率；语言模型对局是探索性观测，而不是具有代表性的行为数据集。因此，恰当的结论必须是条件性的：在所述规则、行动者视角、似然模型、后续策略和随机机制下，不同廉价磋商动作可以具有实质不同的终局价值估计，且这些差异能够由可复现矩阵测量。

下一阶段的目标是经验校准。预注册对局应记录决策时公开状态、行动者合法视角、结构化动作、随后票型和终局结果，用于估计信念反应、比较矩阵事前价值与匹配结果，并单独估计整局处理效应；此后才能评价矩阵作为预测器或决策辅助工具的有效性。

\section{边界以及风险}

\subsection{推断边界}

最重要的边界是蒙特卡洛无偏性的含义。如果轨迹确实从所声明的后验分布、规则随机性和固定功利等级策略中独立生成，且没有按结果删样，则样本均值只对该声明的模拟分布无偏；它不对人类玩家、某个语言模型、未知总体或不同规则集无偏。蒙特卡洛标准误只衡量有限模拟造成的数值噪声，不包含似然模型、策略分数、角色抽象或行为设定错误所造成的不确定性。

第二项边界涉及因果含义。矩阵内配对差在模型中具有清楚的干预语义，因为候选和参照共享相同状态与随机流；探索性整局分组不具有相同控制。每个战术只有2局时，观测胜率差无法与站位、顺序和模型抽样分离，只能被描述为机制观测。未来因果设计需要预先规定随机分组、足够重复、处理依从性检查及处理专属不确定性区间。

第三项边界涉及不同证据来源的可比性。Wang 的理论基线、当前矩阵和较早语言模型语料不是完全相同的游戏。Wang 研究更小角色集，并在特定交流约束下证明结论；当前矩阵加入女巫与守卫并冻结独立结算规则，较早语料则遵循采集时记录的引擎配置。三者可以相互提供研究动机，但在没有共同方案时不能合并数值估计。

\subsection{技术与建模风险}

\begin{longtable}{>{\bfseries\raggedright\arraybackslash}p{0.22\textwidth}>{\raggedright\arraybackslash}p{0.35\textwidth}>{\raggedright\arraybackslash}p{0.34\textwidth}}
\caption{主要风险与计划控制措施}\\
\toprule
风险 & 后果 & 控制措施或下一项检验 \\
\midrule
\endfirsthead
\toprule
风险 & 后果 & 控制措施或下一项检验 \\
\midrule
\endhead
观察者信息泄漏 & 行动者使用其角色依法不可能知道的信息，导致收益查询失效。 & 分离观察者和行动者通道；检验行动者无法区分的隐藏世界，并在决策边界拒绝完整站位。 \\
策略设定错误 & 可复现估计可能描述不合理的续局模型。 & 比较预先声明的替代策略；在留出对局上校准行动频率，且不得追溯修改已报告请求身份。 \\
证据强度校准不足 & 0、0.5和0.8可能不符合人类或模型真实更新方式。 & 只把它们视为敏感性锚点；通过重复的信念引出数据估计反应似然。 \\
小样本排序 & 终局效用离散时，10轨迹案例可能发生顺序反转。 & 正式决策使用每格100轨迹生产预算，同时报告收益与配对差标准误。 \\
多重比较 & 大量目标和证据列使研究者容易过度强调极端行。 & 保留全部行，预先规定核心对照，并用独立种子和对局验证候选动作。 \\
处理不依从 & 语言模型可能不表达已启用战术，或表达被禁止战术。 & 审计结构化动作与完整文本；把偏离如实记录而不是静默重新分类。 \\
规则漂移 & 平票、角色能力或胜负优先级变化会改变估计对象。 & 为每张矩阵附加稳定规则身份；规则不兼容时重新计算而非复用旧结果。 \\
动作抽象损失 & 去除自然语言措辞可能遗漏说服风格、模糊性或语用背景。 & 保留结构语义以支持比较，并把受控语言变体作为独立实验层加入。 \\
递归决策膨胀 & 未来玩家再次构建矩阵会改变目标策略并造成计算增长。 & 固定非递归后续策略；把策略递归作为后续独立模型研究。 \\
均衡过度声明 & 最高估计行可能被误写为全局最优或 PBE。 & 使用“模型条件”限定语；只有另行求解扩展博弈并验证信念与序贯理性后才讨论均衡。 \\
\bottomrule
\end{longtable}

\subsection{下一阶段证据计划}

下一阶段应预注册板子、发言顺序分配、处理因子、语言模型版本、温度、种子方案、重试规则和主要结果。由于两例平民虚假声明均呈现明确反噬机制，经验研究应优先扩充整局对照组和平民虚假声明组。主要票型结果应包括首日放逐预言家、首日放逐狼人、票数集中程度、身份声明竞争频率和预言家存活；终局阵营胜负继续作为延迟结果。夜间战术应作为独立因素研究，而不能混入首日发言效果。

对于矩阵校准，每次发言决策都应在行动发生前记录。记录必须包含公开状态、角色视角私有事实、行动者实际发言位置、结构化动作、后续公开投票和终局结果。观察者真值可以用于评价，但必须保持在动作查询之外。信念反应可以通过引出后续行动者在发言后的角色概率进行测量，也可以在训练集上拟合预先声明的似然模型，再在留出对局上评价。

最后，实质性结论应分阶段提升。稳定矩阵结果需要生产预算重复与不确定性报告；跨模型结果需要在预先声明的不同策略或语言模型配置下重复；行为结论需要留出的观测对局；因果处理结论需要具有充分统计效力的随机整局分组；均衡结论则需要为扩展博弈明确解概念，并提供证明或经过验证的计算。该证据阶梯保持团队已经构建的内容、当前模拟所提示的现象以及未来证据可能支持的结论三者之间的边界。

\begin{thebibliography}{9}

\bibitem{wang2025}
S. Wang，
\emph{Optimal Strategy in the Werewolf Game: A Theoretical Study}，
arXiv:2408.17177v2，2025。

\bibitem{braverman2008}
M. Braverman、O. Etesami 与 E. Mossel，
“Mafia: A theoretical study of players and coalitions in a partial information environment”，
载于 \emph{Proceedings of the 3rd International Conference on Algorithmic Game Theory}，2008，825--846页。

\bibitem{yao2008}
E. Yao，
\emph{A Theoretical Study of Mafia Games}，
arXiv:0804.0071，2008。

\bibitem{migdal2010}
P. Migdal，
\emph{A Mathematical Model of the Mafia Game}，
arXiv:1009.1031，2010。

\bibitem{team2026midterm}
研究团队，
《狼人杀廉价磋商的精确信念收益矩阵中期报告》，
内部研究报告，2026。

\bibitem{team2026experiment}
研究团队，
《廉价磋商研究报告：第一天战术对票型与胜率的影响》，
内部探索性实验报告，2026年8月23日。

\end{thebibliography}

\end{document}
